\chapter{Conclusions and Future Work}
\label{ch:conclusion}

This report has developed, from first principles and in full mathematical detail,
a framework for the distributed real-time estimation and control of heterogeneous
UAV clusters. Two vehicle models were derived completely: a hexacopter with unit-%
quaternion attitude dynamics---Hamilton kinematics, Newton--Euler translational
and rotational dynamics, rotor thrust/torque, and the over-actuated six-rotor
mixer with its rank structure and fault tolerance---and a fixed-wing aircraft in
the Beard--McLain six-degree-of-freedom formulation, with the wind triangle, the
sigmoid-blended stall aerodynamics, the full force and moment coefficient
expansions and the product-of-inertia rotational dynamics, on the fully
parameterised Aerosonde airframe. Every frame, relation and Jacobian was stated
explicitly and validated numerically.

Around a per-agent core of two constrained real-time solvers---a model-predictive
controller and a moving-horizon estimator, both imposing general nonlinear
constraints on the full problem and executing one real-time iteration per
sample---the report constructed a distributed moving-horizon
estimator, in which agents fuse relative measurements to stay localised through
GNSS outages, and a distributed model-predictive controller, coupling the agents
through consensus on a shared reference, decomposition of the coupling
constraints, and control-barrier-function collision avoidance. A three-language
software architecture (C core, C++ coordination, Python research) and a companion
implementation realise the framework. The simulation evidence is statistical,
not anecdotal: beyond the representative missions, a $126$-run Monte-Carlo
campaign over noise realisations, denial severities, degraded graphs and an
estimator ablation recorded no divergences and no solver failures, established
tight percentile bounds on the estimation and formation errors, exposed
\emph{anchor coverage} of the communication graph as the structural condition
for GNSS-denied resilience, and quantified the closed-loop cost of losing the
distributed estimator ($4.8\times$ worse formation keeping and erosion of the
safety margin). A computational study of the embedded solvers instantiated on
the drone models measured microsecond-scale constrained solves---including the
full $13$-state quaternion NMPC---with order-of-magnitude warm-start gains at
identical accuracy; a leveled multi-hop anchoring extension of the DMHE then
repaired the anchor-coverage failures exposed by the campaign; and an
in-the-loop ablation study substituted the embedded C implementations into the
closed loop under paired noise, reproducing the reference behaviour (and
improving the estimation) with zero safety violations. The complete
distributed estimation-and-control loop has thus been validated with the
embedded solver implementations in closed loop on the nonlinear multi-agent
plant---the claim that makes this real-time optimisation machinery a
demonstrated, not merely plausible, basis for multi-agent UAV deployment.

\paragraph{Future work.} Several extensions follow naturally from the structure
built here.
\begin{itemize}[nosep,leftmargin=1.4em]
  \item \textbf{Joint distributed estimation and control.} The estimator and
        controller are a matched pair; coupling their real-time iterations
        directly---so that each agent runs a single distributed estimation-and-%
        control step per sample---is a coherent next step and a genuinely novel
        research direction beyond implementing either alone.
  \item \textbf{Full arrival-cost consensus.} The leveled multi-hop anchoring
        demonstrated here restores coverage whenever some agent retains an
        absolute fix; extending to the full arrival-cost consensus of
        \cite{battistelli2019distributed} would additionally handle the case in
        which \emph{all} agents are intermittently absolute-denied, with
        centralised-like stability.
  \item \textbf{Tighter DMPC coupling.} Writing the collision-avoidance and
        connectivity constraints directly into the MPC horizon (rather than the
        safety filter) and closing the ADMM loop over the network would give
        recursive-feasibility guarantees over the prediction horizon, not only
        instantaneously.
  \item \textbf{Hardware-in-the-loop and flight.} Running the C core on a
        companion computer against a PX4 software-in-the-loop stack, and then in
        flight, would validate the timing and robustness claims on real hardware
        and real communication links.
  \item \textbf{Heterogeneous missions.} Mixed hexacopter/fixed-wing clusters---%
        loitering fixed-wing relays providing communication and GNSS anchoring to
        manoeuvring multirotors---are directly supported by the interface-based
        design and are a compelling application of the heterogeneity the
        framework was built to accommodate.
\end{itemize}
